\item \textbf{A dissection of the audit's two surviving positives.} Placing quantum and classical kernels
on identical data with exact Gram matrices and equal tuning, we show that a quantum-kernel advantage
surviving both FDR and Holm correction is produced by the classical baseline's bandwidth search, because
in-distribution cross-validation cannot see NSL-KDD's train-to-test shift; controls on the test set
alone show that cross-validation is reliable without the shift and apportion the artefact between novel
attack types and a changed attack mix. For the audit's other positive, a four-qubit hybrid's advantage at
$1\%$ false positives, we build classical twins that replace only the quantum layer, show that the
trained circuit is exactly an 81-term trigonometric polynomial, and find that a tanh layer of the same
width reproduces the result, whose margin over the Random Forest lies entirely in novel attack types.
\item \textbf{An induced-shift protocol across three datasets and a label-free warning statistic.}
Holding attack families out of training on UNSW-NB15, ToN-IoT and CICIDS2017 (26 runs) recreates the
selection failure, including a spurious $0.17$ quantum lead on CICIDS2017, while at its best setting the
quantum kernel never leads by more than $0.003$. The squared MMD between training and unlabelled test
traffic tracks how badly cross-validation fails. Leave-attack-types-out cross-validation repairs
selection on NSL-KDD but not under induced shift, so we recommend reporting kernel comparisons under
several selection rules with the oracle bound. The geometric difference of Huang et al., on all four
datasets, shows the IQP kernel to be reproducible by a classical kernel on its own circuit phases.
\item \textbf{Hardware validation of both models.} The projected kernel and the hybrid's quantum layer
run on three IBM Heron processors: the kernel's classifier stays within $0.01$ ROC-AUC of simulation,
and the hybrid's low-false-positive detection is unchanged on 6{,}000 test rows per device and
indistinguishable from its classical twin's on the same rows.
